\subsection*{Общая характеристика работы}

\newcommand{\actuality}{\underline{\textbf{Актуальность темы.}}}
\newcommand{\aim}{\underline{\textbf{Целью}}}
\newcommand{\tasks}{\underline{\textbf{задачи}}}
\newcommand{\defpositions}{\underline{\textbf{Основные положения, выносимые на~защиту:}}}
\newcommand{\novelty}{\underline{\textbf{Научная новизна:}}}
\newcommand{\influence}{\underline{\textbf{Практическая значимость}}}
\newcommand{\reliability}{\underline{\textbf{Достоверность}}}
\newcommand{\probation}{\underline{\textbf{Апробация работы.}}}
\newcommand{\contribution}{\underline{\textbf{Личный вклад.}}}
\newcommand{\publications}{\underline{\textbf{Публикации.}}}

{\actuality}
 Актуальность работы обусловлена планами создания пилотируемого транспортного корабля нового поколения, работающего на высокоширотных и лунных орбитах. Проект транспортного корабля активно разрабатывается с 2010~г. и к настоящему времени, начата работа по выпуску рабочей конструкторской документации на составные части корабля, в том числе и на дозиметр бортовой.
 
 Несмотря на непрерывный дозиметрический контроль всех российских космических миссий, начиная с первого полета человека в космос и заканчивая полетами экспедиций на МКС, не вызывает сомнений необходимость продолжения ряда исследований радиационной обстановки на каждом из пилотируемых и на значительной части беспилотных космических аппаратах.  
 
 Именно поэтому необходимо разработать приборы  для проведения непрерывного дозиметрического мониторинга области околоземного пространства, в которой планируется проведение перспективных пилотируемых полетов. \textit{Данная работа направлена на создание основ для осуществления такого мониторинга.}

 \aim\ данной работы является разработка методов исследования распределения мощности дозы космической радиации и создание на основе этих методов современных приборов, предназначенных для космических аппаратов работающих на околоземных и лунных орбитах.
  

Для~достижения поставленной цели в работе необходимо было решить следующие задачи {\tasks}:
\begin{enumerate}
  \item Систематизация и обобщение характеристик радиационных условий на аналогичных орбитах (аппараты БИОН, Прогноз, Cluster, POES) для разработки программы эксперимента;
  \item Разработать требования к бортовому дозиметру для нового пилотируемого транспортного корабля;
  \item Разработать прибор для дозиметрического мониторинга на борту космического аппарата <<Ломоносов>>;
  \item Участвовать в подготовке и проведении эксперимента с дозиметром на борту КА <<Ломоносов>>;
  \item Обработать и провести анализ научной информации с прибора ДЭПРОН.
  
\end{enumerate}

\defpositions
\begin{enumerate}
  \item Разработан прибор для радиационного мониторинга на борту КА
  \item Подтверждены измерительные характеристики нового прибора
  \item Получены и обработаны дозиметрические экспериментальные данные
  \item Выделен вклад в дозы при пересечении различных областей космического пространства - внутреннего и внешнего радиационного пояса
\end{enumerate}

\novelty
\begin{enumerate}
  \item Впервые разработан исследовательский прибор, сочетающий в едином блоке дозиметр заряженных частиц и нейтронные детекторы;
  \item Было выполнено оригинальное моделирование спектрометрических свойств дозиметра протонов электронов и нейтронов.
  \item Было выполнено создание базы данных спутниковых измерений прибора Дэпрон за все время работы прибора.
\end{enumerate}

\influence\ определяется необходимостью контроля радиационных условий для целей обеспечения надежной работы аппаратуры на полярных спутниках. Прибор Дэпрон позволил провести полноценной исследование всех основных компонент радиационного излечения, вносящих вклад в поглощенную и эквивалентную дозы на борту космического аппарата. 

\reliability\ полученных результатов обеспечивается публичностью полученных данных спутниковых измерений радиационных условий. Результаты находятся в соответствии с результатами, полученными другими авторами приборах RELEC и ELFIN, а также на других полярных спутниках (POES). Результаты измерений сравнивались с измерениями на баллонах BARREL (запуски август 2016г).

\probation\
Основные результаты работы докладывались~на:
Ломоносовских чтениях, международной конференции COSPAR2014, конференции ОМУС-2015 (г. Дубна), нескольких конференциях Ломоносовские чтения, рабочих совещаний коллаборации Lomonosov, а также ряде других конференций.

\contribution\ Автор принимал активное участие в проведении тепловых испытаний полупроводникового детектора с усилительным трактом прибора. Проведены проверки работоспособности усилительного тракта с радиоактивными источниками ОСГИ Cs137 и Co60. Автор участвовал в проведении испытаний детекторов тепловых нейтронов на лабораторном источнике нейтронов. 

Проведены работы по стыковке и согласованию платы цифровой обработки сигналов с аналоговыми усилительными трактами и дискриминирующими блоками прибора. Автором вместе с его научным руководителем написана программа на С++ для контроллера платы цифровой обработки сигналов. Для наземной отработки и испытаний написана программа для ПК на WinForms/С\#, позволяющая оперативно контролировать параметры работы прибора и выходные данные. 

Создана исчерпывающая модель дозиметра в системе Catia и подготовлена для использования в Монте-Карло моделировании. Написана программа на базе пакета Geant4 для математического моделирования характеристик прибора Дэпрон, а также программы для моделирования характеристик приборов ДБ-8 и других приборов для радиационных измерений. 

Атором подготовлены программы для анализа данных прибора Дэпрон, в том числе программы предобработки бинарных данных и программы визуализации данных детекторов прибора. 
Автором подготовлена методика точной привязки спутниковых данных с прибора Дэпрон ко всемирному времени. С использованием оригинальной методики получена полная база данных измерений прибора Дэпрон за все время работы прибора.

\publications\ Основные результаты по теме диссертации изложены в ХХ печатных изданиях~\cite{zolotarev2017numerical51590279,vlasova2017optimization27547274,vlasova2015operational11246447},
Х из которых изданы в журналах, рекомендованных ВАК~\cite{zolotarev2016chislennoe32150868}, 
ХХ --- в тезисах докладов~\cite{zolotarev2016modelirovanie36997161,grafodatsky2016development21020642, amelyushkin2015sozdanie1065732, angelopoulos2010university5851829, angelopoulos2011university1295357}.
 % Характеристика работы по структуре во введении и в автореферате не отличается (ГОСТ Р 7.0.11, пункты 5.3.1 и 9.2.1), потому её загружаем из одного и того же внешнего файла, предварительно задав форму выделения некоторым параметрам

%Диссертационная работа была выполнена при поддержке грантов ...

%\underline{\textbf{Объем и структура работы.}} Диссертация состоит из~введения, четырех глав, заключения и~приложения. Полный объем диссертации \textbf{ХХХ}~страниц текста с~\textbf{ХХ}~рисунками и~5~таблицами. Список литературы содержит \textbf{ХХX}~наименование.

%\newpage
\subsection*{Содержание работы}
Во \underline{\textbf{введении}} обосновывается актуальность исследований, проводимых в рамках данной диссертационной работы, приводится обзор научной литературы по изучаемой проблеме, формулируется цель, ставятся задачи работы, сформулированы научная новизна и практическая значимость представляемой работы.

\underline{\textbf{Первая глава}} посвящена ...

 картинку можно добавить так:
\begin{figure}[ht] 
  \center
  \includegraphics [scale=0.27] {latex}
  \caption{Подпись к картинке.} 
  \label{img:latex}
\end{figure}

Формулы в строку без номера добавляются так:
\[ 
  \lambda_{T_s} = K_x\frac{d{x}}{d{T_s}}, \qquad
  \lambda_{q_s} = K_x\frac{d{x}}{d{q_s}},
\]

\underline{\textbf{Вторая глава}} посвящена исследованию 

\underline{\textbf{Третья глава}} посвящена исследованию 

В \underline{\textbf{четвертой главе}} приведено описание 

В \underline{\textbf{заключении}} приведены основные результаты работы, которые заключаются в следующем:
%% Согласно ГОСТ Р 7.0.11-2011:
%% 5.3.3 В заключении диссертации излагают итоги выполненного исследования, рекомендации, перспективы дальнейшей разработки темы.
%% 9.2.3 В заключении автореферата диссертации излагают итоги данного исследования, рекомендации и перспективы дальнейшей разработки темы.

В процессе разработки аппаратуры для радиационного мониторинга на борту космического аппарата были выделены основные требования к дозиметру, которые позволили создать прибор ДЭПРОН, удовлетворяющий и выполнению задачи мониторинга обстановки и исследовательским целям в космической дозиметрии. Для решения этой задачи был проведен анализ современной литературы по направлению космической дозиметрии. Были найдены аналоги создаваемого прибора и определен ряд критических параметров для такого типа приборов, в том числе энергетические диапазоны регистрируемых излучений, типы излучений и ориентировочные потоки в изучаемых областях космического пространства. Эти параметры и определили оптимальные размеры детекторов, их расположение и требуемое быстродействие встраиваемых вычислительных процессоров.

Для выполнения поставленных задач был создан активный дозиметр нового типа с возможностью регистрации нейтронов тепловых энергий.

Средняя доза на 

\section{Выводы}
\begin{enumerate}
	\item Систематизированы и обобщены характеристики радиационных условий на аналогичных орбитах (аппараты БИОН, Прогноз, Cluster, POES) для разработки программы эксперимента;
	\item Разработаны требования к бортовому дозиметру для нового пилотируемого транспортного корабля;
	\item Разработан прибор для дозиметрического мониторинга на борту космического аппарата <<Ломоносов>>;
	\item Подготовлен и проведен эксперимент с дозиметром на борту КА <<Ломоносов>>;
	\item Обработана полученная с прибора ДЭПРОН информация и проведен её анализ.	
\end{enumerate}

Дополнительные результаты, полученные в результате проведения эксперимента ДЭПРОН: 
\begin{enumerate}

  \item Приближенные оценки показали, что полупроводниковые детекторы прибора чувствительны к электронам энергий более 0,5 МэВ и протонам с энергиями более 5 МэВ;
  \item Математическое моделирование показало что максимум функции чувствительности нейтронный счетчиков соответствует энергии нейтронов 0,005 МэВ и 0,05 МэВ детекторов различной защищенности;
  \item Поглощенная доза за время проведения эксперимента достигла ... и ... для верхнего и нижнего детекторов.
  
\end{enumerate}



%\newpage
\renewcommand{\refname}{\large Публикации автора по теме диссертации}
%\nocite{*}
%\insertbiblioauthor                          % Подключаем Bib-базы
\insertbibliofull
